\documentclass[11pt]{article}
\usepackage[a4paper,margin=27mm]{geometry}
\usepackage{amsmath,amssymb,amsthm}
\usepackage{array}
\usepackage[T1]{fontenc}
\usepackage{hyperref}
\hypersetup{colorlinks=true,linkcolor=blue,urlcolor=blue,
 pdftitle={A cofinal bound for maximum-modulus points of entire functions},
 pdfauthor={Yunpeng Li}}
\newtheorem{theorem}{Theorem}
\newtheorem{lemma}[theorem]{Lemma}
\newtheorem{proposition}[theorem]{Proposition}
\newtheorem{corollary}[theorem]{Corollary}
\numberwithin{equation}{section}
\newcommand{\C}{\mathbb C}
\newcommand{\N}{\mathbb N}
\newcommand{\R}{\mathbb R}
\newcommand{\Cs}{\C^{\times}}
\newcommand{\Chat}{\widehat{\C}}
\newcommand{\Max}{\operatorname{MaxPoints}}
\newcommand{\card}{\operatorname{card}}
\newcommand{\ord}{\operatorname{ord}}
\newcommand{\cl}{\operatorname{cl}}
\title{A cofinal bound for maximum-modulus points of entire functions}
\author{Yunpeng Li\\[3pt]
 \small School of Data Science, The Chinese University of Hong Kong, Shenzhen, China\\
 \small\texttt{liyunpeng@cuhk.edu.cn}}
\date{11 October 2026}
\begin{document}
\maketitle

\begin{abstract}
For every nonzero entire function which is not a monomial, we prove that
the number of its exact maximum-modulus points is bounded by a fixed finite
integer at arbitrarily large radii. This gives a negative answer to the
all-radii maximum-modulus question. After removing the zero at the origin,
one may take the bound to be twice the order at zero of the normalized
logarithmic derivative minus its constant term. The proof uses finite local
inverse-root families and continuation of the closure of regular targets
with large fibers. It handles the omitted logarithmic-derivative value and
the value at infinity without assuming global normalization or a generic
component degree. A complete Lean~4 formalization accompanies the proof.
\end{abstract}

\section{The problem and the result}\label{sec:introduction}

Throughout, $\N=\{0,1,2,\ldots\}$, so the monomials excluded below include
nonzero constant functions.

For an entire function $f$ and a positive radius $r$, write
\[
 M_f(r)=\max_{|z|=r}|f(z)|.
\]
\[
 \Max(f,r)=\{z\in\C:|z|=r\text{ and }
             |f(w)|\le |f(z)|\text{ for every }w\text{ with }|w|=r\}.
\]
Thus the maximum-point set is defined by actual function values. The maximum
exists by compactness. For a nonzero non-monomial entire function, we prove
below that this set is finite at every positive radius; its cardinality
$\nu_f(r)$ consequently has its ordinary finite meaning.

Hayman--Lingham's Problem~2.16(b), on printed page~29 of
\cite{HL}, asks whether a single entire function can have
$\liminf_{r\to\infty}\nu_f(r)=\infty$. The non-monomial formulation is
explicit in Gl\"ucksam--Pardo-Sim\'on's Question~1.1(b) \cite{GP} and in
the introduction of Pardo-Sim\'on--Sixsmith \cite{PS}. In contrast,
Herzog--Piranian \cite{HP} answered the selected-radii question
$\limsup_{r\to\infty}\nu_f(r)=\infty$ affirmatively in 1968. These two
questions have different quantifiers. The recent finite-order result
\cite{PS} concerns the latter question. Approximate maximum points,
as studied in \cite{GP}, also differ from the exact set used here.

\begin{theorem}\label{thm:main}
Let $f$ be a nonzero entire function which is not $cz^m$ for any
$c\ne0$ and $m\in\N$. There is an integer $B\in\N$ such that, for every
$R>0$, there is an $r>R$ for which $\Max(f,r)$ is finite and
$\card\Max(f,r)\le B$.
In fact, write $f(z)=z^m g(z)$ with $g(0)\ne0$. One may take
\[
 B=2k,\qquad k=\ord_0\left(\frac{zg'(z)}{g(z)}\right).
\]
Here $k$ is a positive finite integer.
\end{theorem}

\begin{corollary}\label{cor:negative}
No nonzero non-monomial entire function has
$\nu_f(r)\to\infty$ as $r\to\infty$ along all positive radii.
\end{corollary}

\begin{proof}[Deduction from Theorem~\ref{thm:main}]
Fix the bound $B$ supplied by the theorem. If the count tended to infinity,
there would be a threshold $R_0$ beyond which it was at least $B+1$.
Choose $R>\max\{0,R_0\}$ and then the radius $r>R$ supplied by the theorem.
Its maximum-point set is finite and has cardinality at most $B$, a
contradiction. The same argument rules out a formulation which treats an
infinite maximum-point set as having infinite cardinality.
\end{proof}

The proof first bounds fibers with a small coordinate. It then constructs
finite local covers of the actual logarithmic-derivative correspondence,
counting physical pairs rather than inverse-root indices. The closure of
the regular targets with too many pairs continues across every spherical
value, including an omitted value. Its product projection would be both
open and closed in the plane, contradicting the small-coordinate bound.
A countable exceptional set and the radial maximum envelope then give a
good radius in every positive interval.

\section{Factorization and finiteness on each circle}\label{sec:finiteness}

Factor $f(z)=z^m g(z)$, where $m\in\N$, $g$ is entire, and $g(0)\ne0$.
The factor $g$ is nonconstant, since otherwise $f$ is a monomial. Put
\[
 g^\sharp(z)=\overline{g(\bar z)}.
\]
This reflection is entire. On every positive circle,
$|f(z)|=r^m|g(z)|$, so $f$ and $g$ have the same maximum-point set.

\begin{lemma}\label{lem:product}
For each $s\in\Cs$, the function
$P_s(z)=g(z)g^\sharp(s/z)$ is holomorphic and nonconstant on $\Cs$.
\end{lemma}

\begin{proof}
Holomorphy follows by composition and multiplication. If $P_s$ were a
constant $C$, then, as $z\to\infty$, the factor $g^\sharp(s/z)$ would tend
to $g^\sharp(0)\ne0$. Outside a sufficiently large disk it would be
bounded away from zero, and the identity
$g(z)=C/g^\sharp(s/z)$ would bound $g$. Continuity bounds $g$ on the
remaining compact disk. Liouville's theorem would make $g$ constant,
contrary to its choice.
\end{proof}

\begin{lemma}\label{lem:circle-finite}
The set $\Max(f,r)$ is finite for every $r>0$.
\end{lemma}

\begin{proof}
It suffices to consider $g$. If its maximum-point set on the compact
circle were infinite, it would have an accumulation point on that circle,
hence in $\Cs$. For every maximum point $z$, the identity
$r^2/z=\bar z$ gives
\[
 P_{r^2}(z)=g(z)g^\sharp(\bar z)=|g(z)|^2=M_g(r)^2.
\]
The identity theorem on the connected domain $\Cs$ would make
$P_{r^2}$ constant. This contradicts Lemma~\ref{lem:product}.
\end{proof}

This proves finiteness independently of the subsequent fiber bound and
justifies using $\nu_f(r)=\card\Max(f,r)$ throughout the positive domain.

\section{Logarithmic derivatives and the small-coordinate bound}
\label{sec:small}

Define
\[
 A(z)=m+\frac{zg'(z)}{g(z)},\qquad B(z)=A^\sharp(z).
\]
These are meromorphic functions interpreted as continuous functions into
$\Chat$. The point zero is regular and $A(0)=B(0)=m$. At a zero
$\alpha\ne0$ of $g$ of multiplicity $n$, write
$g(z)=(z-\alpha)^n h(z)$ with $h(\alpha)\ne0$. Then
\[
 A(z)=\frac{n\alpha}{z-\alpha}+m+n+\frac{zh'(z)}{h(z)}.
\]
Thus every pole is genuine and simple, and its spherical value is
$\infty$. The same holds for $B$.

The function $A-m$ is analytic near zero and not identically zero. Indeed,
if it vanished on a neighborhood, $g'$ would vanish on a punctured
neighborhood where $g\ne0$, and the identity theorem would make $g$
constant. Let
\[
 k=\ord_0(A-m)>0,\qquad L=2k.
\]
Reflection gives the same positive finite order for $B-m$.

For any product $s$ and spherical common value $p$, let
\[
 F(s,p)=\{(z,w)\in\C^2:zw=s,\ A(z)=B(w)=p\}.
\]
We count distinct actual ordered pairs.

We will repeatedly use the following local power coordinate. If a
nonconstant analytic function $h$ satisfies $h(\alpha)=p_0$, write
\[
 h(z)-p_0=(z-\alpha)^d u(z),\qquad u(\alpha)\ne0.
\]
After shrinking to a simply connected neighborhood, $u$ has an analytic
$d$th root. The coordinate $e(z)=(z-\alpha)u(z)^{1/d}$ has
$e'(\alpha)\ne0$, so the inverse function theorem makes it conformal and
$h=p_0+e^d$. In that coordinate, each value has at most $d$ distinct
inverse roots. Apply this to $A-m$ and $B-m$ at zero. Choose
$\varepsilon>0$ such that both coordinate disks have this $k$-root bound
and contain no poles.

For $s\ne0$, a pair in $F(s,p)$ is determined by either coordinate, since
$w=s/z$ and $z=s/w$. Thus at most $k$ pairs have $|z|<\varepsilon$, and
at most $k$ pairs have $|w|<\varepsilon$. Their union has at most $L$ pairs,
also when $p=\infty$. In particular,
\begin{equation}\label{eq:small}
 0<|s|<\varepsilon^2\quad\Longrightarrow\quad
 F(s,p)\text{ is finite and }\card F(s,p)\le L
 \quad\text{for every }p\in\Chat.
\end{equation}
Indeed, a pair whose two coordinate norms were at least $\varepsilon$
would have $|zw|\ge\varepsilon^2$.

\section{Properness and exact finite local covers}\label{sec:covers}

Consider the actual correspondence
\[
 X=\{(z,w)\in\Cs\times\Cs:A(z)=B(w)\ne m\}
\]
over
\[
 \Omega=\Cs\times(\Chat\setminus\{m\}),
\]
and the map $\Phi:X\to\Omega$, $\Phi(z,w)=(zw,A(z))$.
It is proper with finite fibers. Indeed, on a compact target set the common
values stay away from $m$. Continuity of $A,B$ at zero gives $\eta>0$
bounding both coordinate norms below. A bound $|zw|\le S$ then bounds
both norms above by $S/\eta$. The inverse image is a closed subset of the
resulting compact annulus product. For a fixed finite value the inverse
roots are discrete, by nonconstancy and the identity theorem; at infinity
they are isolated poles. Hence there are finitely many in the compact
annulus, and the product fixes the other coordinate. Denote the actual
target image in $\Omega$ by $Y$. Properness makes $Y$ closed in $\Omega$.

All local target neighborhoods in this section lie inside $\Omega$.
Fix $y_0=(s_0,p_0)\in\Omega$. Its actual central fiber is finite.
If it is empty, properness gives a target neighborhood disjoint from $Y$,
so the local cover there is empty. Otherwise, properness captures every
nearby pair in a finite collection of chart products around the central
pairs. To see this, an uncaptured sequence over targets tending to $y_0$
would have a subsequence converging to a central pair inside one of those
chart products.

In a finite-value chart, the analytic power-coordinate construction gives
\[
 A(z)=p_0+e_A(z)^{d_A},\qquad
 B(w)=p_0+e_B(w)^{d_B},
\]
where $e_A,e_B$ are conformal coordinates centered at their physical roots.
Choose one positive common multiple $N$ of the finitely many orders.
For $p=p_0+t^N$, every local inverse root is of the form
\[
 z=e_A^{-1}(\zeta t^{N/d_A}),\qquad
 w=e_B^{-1}(\xi t^{N/d_B}),\qquad
 \zeta^{d_A}=\xi^{d_B}=1.
\]
This constructs a finite analytic physical family $P_i(t)=(z_i(t),w_i(t))$,
with product functions $S_i(t)=z_i(t)w_i(t)$ and raw image maps
\begin{equation}\label{eq:raw}
 \gamma_i(t)=(S_i(t),p_0+t^N).
\end{equation}
All displayed pairs have the actual common values. Every nearby actual
pair is represented by an index at that same $t$, because the roots of each
power equation are exactly the displayed root-of-unity multiples. At zero
every central physical pair is represented. Indices need not be injective.

Localize this capture in the target as follows. Choose a product disk about
$s_0$ and a value disk whose product lies in the capture neighborhood.
Choose $\delta>0$ small enough that all raw product images on $|t|<\delta$
lie in that product disk and $\delta^N$ is smaller than the initial value
radius. Use the value radius $\delta^N$ in the final neighborhood $W$.
The full raw open-disk images then lie in $W$, and every target in $W$
has all its parameter roots inside that same disk. Consequently
\begin{equation}\label{eq:cover}
 Y\cap W=\bigcup_i\gamma_i(\{|t|<\delta\}).
\end{equation}
This is an exact cover of the total actual image.

At $p_0=\infty$, use $\lambda=1/p$. Every pole of $A$ and $B$ is simple,
so $1/A$ and $1/B$ have simple zeros at the corresponding physical roots.
They are conformal coordinates, and the physical roots are analytic
functions of $t=\lambda$ itself. Here take $N=1$ and $p(t)=1/t$ for
$t\ne0$, with $p(0)=\infty$. Properness and the same target shrinking give
the exact cover. Its punctured points have finite common value. In either
chart type, if $v$ is the centered value coordinate, then
$|v(p(t))|=|t|^N$.

\section{Stabilized branches and distinct physical counts}
\label{sec:stabilization}

For finitely many indices $i,j$ and $N$th roots of unity $\omega$, the
functions
\[
 S_i(\omega t)-S_j(t)
\]
are analytic at zero. By the isolated-zero theorem, after a common shrinking
each is either identically zero as a germ or nonzero for every nonzero
parameter. Stabilize the comparisons of both physical coordinates as well.
Repeat the matched target shrinking above after this change of radius.
These finite comparisons give the following facts.

\begin{enumerate}
\item Two punctured raw image representatives are equal or disjoint.
An intersection has parameter roots related by $t\mapsto\omega t$;
the comparison is therefore a persistent identity. Rotation preserves
the parameter disk.

\item Every noncentral point of the finite raw image union has a
neighborhood where the whole union is one analytic graph over the common
value coordinate. The map $t\mapsto t^N$ has nonzero derivative there.
Coincident rotated roots give the same graph germ by the persistent
comparison. Noncoincident roots and the other disjoint branches are absent
after shrinking the neighborhood. Relative closedness, proved below,
also excludes limits from those other branches.

\item Each punctured image branch has a constant physical fiber cardinal
$K$. Exact capture at the same parameter gives, for each index $j$,
\begin{equation}\label{eq:exact-fiber}
 F(S_j(t),p(t))=\{P_i(t):S_i(t)=S_j(t)\}.
\end{equation}
The product filter and the equivalence relation given by physical equality
$P_i(t)=P_{i'}(t)$ are constant on the punctured disk, by the stabilized
comparisons. This finite image therefore has constant cardinality. It
counts distinct pairs even when several indices represent the same pair.

\item The finite raw image union, and any subunion of its branches, is
relatively closed in its open value cylinder. If targets converge inside
that cylinder, $|v(p)|=|t|^N$ bounds their parameters strictly inside the
disk. Passing to a fixed index and then to a convergent parameter
subsequence supplies the target limit, including at the center.
\end{enumerate}

There is an essential central-cardinality comparison. Suppose the entire
actual image at a finite center is locally a single graph $s=\psi(p)$.
Forward validity forces every product germ to be $\psi(p_0+t^N)$, so all
indices belong to one persistent product class. Choose index representatives
of distinct central physical pairs. If two representatives coincide at a
small nonzero parameter, stabilized physical equality makes their germs
equal, and hence their central pairs equal. The distinct chosen central
representatives thus remain distinct at that parameter. Therefore
\begin{equation}\label{eq:centralcount}
 \card F(s_0,p_0)\le K.
\end{equation}
This uses the graph property of the total image and the exact cover.

\section{The regular high-count closure}\label{sec:high-count}

Call a point $(s_0,p_0)\in Y$ \emph{regular} if in some product
neighborhood $U_s\times U_p$ the total image satisfies
\[
 Y\cap(U_s\times U_p)=\{(\psi(p),p):p\in U_p\}
\]
for an analytic function $\psi$ with values in $U_s$. The value coordinate
is spherical when needed. This is equality for the whole local image.

Let $E$ be the set of regular finite-valued targets in $Y$ whose actual
fiber has more than $L$ distinct pairs, and let
\[
 H=\cl E\quad\text{in }\C\times\Chat.
\]
By \eqref{eq:small}, $E$ misses the full open cylinder
$\{|s|<\varepsilon^2\}\times\Chat$: its first coordinate is nonzero by
definition, and every nonzero product in that cylinder has at most $L$
pairs. The same open cylinder is disjoint from $H$. In particular, every
point of $H$ has nonzero product coordinate. Closedness of $Y$ in $\Omega$
also gives $H\subset Y$ locally away from $p=m$.

In a stabilized local finite family, select the branches whose punctured
physical count $K$ exceeds $L$. On the matched neighborhood $W$,
\[
 \text{selected punctured union}\ \subset\ E\cap W\
 \subset\ \text{selected full union}.
\]
Away from the center this follows from the exact whole-image graph and
count. A finite central point of $E$ belongs to a selected branch by
\eqref{eq:centralcount}; it cannot contribute an isolated point.
At infinity the center itself is not in $E$, but the punctured pole-chart
points are finite-valued and satisfy the same selection rule.
Relative closedness and density of each punctured branch in its full
image give
\begin{equation}\label{eq:interiorH}
 H\cap W=\text{selected full union}\qquad(p_0\ne m).
\end{equation}
Closure is localized inside the open neighborhood: every sequence tending
to a point of $W$ is eventually in $W$. Each point of $H$ away from $m$
therefore lies on a complete local analytic branch contained in $H$,
including when its value is infinity.

\section{Continuation across the omitted value}\label{sec:omitted}

Fix $y_0=(s_0,m)\in H$; the preceding exclusion gives $s_0\ne0$.
Choose a target neighborhood with $|s|\le S$.
A point of $E$ there has more than $L$ pairs, of which at most $L$
have a small coordinate. At least one surviving actual pair satisfies
\begin{equation}\label{eq:survivor}
 \varepsilon\le |z|,|w|\le S/\varepsilon.
\end{equation}
The upper bounds follow from the product bound and the lower bounds.
The surviving pair can change from target to target; no persistent choice
of it is needed.

The roots of $A-m$ and $B-m$ in this compact annulus are finite.
Finite-value power charts at these roots construct a finite common-power
family across $p=m$. Compactness captures all pairs satisfying
\eqref{eq:survivor} in the chosen chart neighborhoods when $p$ is
sufficiently close to $m$: a contrary sequence would have a subsequence
converging to a captured central root. Hence every nearby point of $E$
lies in this finite raw image union $Z$. This family need not enumerate
escaping pairs. One bounded surviving pair per target gives the containment.

Discard germs whose central product differs from $s_0$. If there are any,
let $d>0$ be the minimum of their finitely many distances $|S_i(0)-s_0|$.
Shrink the parameter disk so $|S_i(t)-S_i(0)|<d/3$ and restrict the target
product to $|s-s_0|<d/3$. None of those germs can represent a target in
that product disk. If there are no discarded germs, the step is vacuous.
Every retained branch has center $y_0$.

Stabilize these finitely many product comparisons as before, with parameter
radius $\delta_0$. Shrink the capture neighborhood $U$ into
$\{|p-m|<\delta_0^N\}$. Then
\[
 E\cap U\subset Z_0,\qquad
 Z_0=\bigcup_j\gamma_j(\{|t|<\delta_0\}),
\]
because every represented parameter over $U$ has $|t|<\delta_0$.
Choose $0<\delta<\delta_0$ so every retained full raw $\delta$-image lies
in $U$, and put
\[
 W=U\cap\{|p-m|<\delta^N\},\qquad
 Z_j=\gamma_j(\{|t|<\delta\}).
\]
Each $Z_j$ lies in $W$. Conversely, a target in $W$ represented in an old
$\delta_0$-image has $|t|^N<\delta^N$, hence $|t|<\delta$.
Using relative closedness in the old value cylinder and the localized
containment $E\cap U\subset Z_0$, we obtain
\begin{equation}\label{eq:boundarycontainment}
 H\cap W\subset\bigcup_j Z_j,\qquad Z_j\subset W.
\end{equation}
Indeed, a sequence from $E$ converging to a point of $W$ is eventually in
$U$, and its represented parameters have a subsequence strictly inside
the old disk. The limit's value then places its parameter inside the
new disk. Matching the value radius with $\delta^N$ is essential to
retain each whole raw $\delta$-image inside $W$.

Write $Z_j^*=\gamma_j(\{0<|t|<\delta\})$.
These punctured images are connected, pairwise equal or disjoint,
and locally smooth graphs over the value. They avoid $p=m$.
Their intersections with $H$ are closed relative to $Z_j^*$.
They are also open, as follows.

At a point $y\in H\cap Z_j^*$, \eqref{eq:interiorH} supplies a complete
analytic subbranch through $y$ contained in $H$. Shrink it inside $W$.
By \eqref{eq:boundarycontainment} and the local single-graph description
of the finite $Z$ union, it is contained in the graph of $Z_j^*$ near $y$.
Its value coordinate is nonconstant. The open mapping theorem makes its
value image a neighborhood of the value at $y$; the graph has one product
coordinate per value, so the subbranch fills a neighborhood of $y$ in
$Z_j^*$. This proves openness of membership in $H$.

Connectedness now gives $Z_j^*\subset H$ or $Z_j^*\cap H=\varnothing$.
All representatives lie inside $W$, so the capture remains valid along
the whole connected punctured image. Since $y_0\in\cl E$ and $E$ contains
no point of value $m$, at least one of the finitely many punctured
branches meets $E$ arbitrarily close to $y_0$. Its punctured image is
contained in $H$, and closedness includes its center.
This selection uses membership in $H$, not a count attached to the
partial bounded-survivor family. Every omitted-value point of $H$
therefore lies on a complete branch contained in $H$; no isolated center
is added.

\section{Nonconstant product coordinates and the global contradiction}
\label{sec:global}

No actual image germ used above has a nonzero constant product.
Suppose $z(t)w(t)=s_0\ne0$ on such a germ. The common-value coordinate
makes $z(t)$ nonconstant: a constant physical coordinate would have
constant common value, whereas the value parameter is nonconstant.
Choose a noncentral parameter with finite common value. Both physical
coordinates then avoid zeros of the respective entire factors.
The open mapping theorem applied to $z(t)$ gives
\[
 A(z)=B(s_0/z)
\]
on an open set where both factors are nonzero. There,
\[
 \frac{zP_{s_0}'(z)}{P_{s_0}(z)}
 =\frac{zg'(z)}{g(z)}
  -\frac{(s_0/z)(g^\sharp)'(s_0/z)}{g^\sharp(s_0/z)}
 =A(z)-B(s_0/z)=0.
\]
The identity theorem makes $P_{s_0}'$ vanish throughout $\Cs$.
This contradicts Lemma~\ref{lem:product}. For a germ centered at poles,
its noncentral common values are finite, so the same argument applies.
No regularity of the total image at that parameter is required.

The complete branches in $H$ at ordinary finite, omitted, and infinite
values therefore have nonconstant analytic product coordinates.
The open mapping theorem makes the product projection of $H$ open,
including at branch centers. The projection is closed because $H$ is
closed in $\C\times\Chat$ and $\Chat$ is compact.
If $E$ were nonempty, this projection would be a nonempty open-and-closed
subset of the connected plane, hence the whole plane.
But the cylinder $|s|<\varepsilon^2$ is disjoint from $H$.
This contradiction proves the following assertion.

\begin{proposition}\label{prop:empty}
Every regular finite-valued target in $Y$ has at most $L$ distinct
physical pairs in its fiber; equivalently, $E=\varnothing$.
\end{proposition}

\section{Countable exceptions and arbitrarily large good radii}
\label{sec:radii}

Let $D$ consist of the nonregular actual image targets in
$\Cs\times(\C\setminus\{m\})$. The finite local family shows that $D$ is
locally finite: all noncentral points in a small local image cover are
single-graph points. At an empty central fiber, properness supplies a
neighborhood outside the image. The ambient space is second countable;
a countable cover by neighborhoods meeting $D$ finitely shows that $D$
is countable. Its product projection is countable, although that projection
need not be locally finite.

Let $Q$ be the positive radii $r$ for which $(r^2,p)\in D$ for some finite
$p$. This is countable since $D$ is countable and $r\mapsto r^2$ is
injective on positive radii. For $r>0$ outside $Q$ and any real $p>m$,
a nonempty fiber $F(r^2,p)$ lies at a regular finite target.
Properness proves it finite, and Proposition~\ref{prop:empty} bounds
its cardinality by $L$. An empty fiber is finite with cardinality zero.
No convention for the cardinality of an infinite set is used.

Set $G(r)=\max_{|z|=r}|g(z)|^2$ for $r>0$.
The maximum principle makes $G$ strictly increasing: equality at two
radii would put a maximum point from the smaller circle in the interior
of the larger disk, forcing $g$ to be constant.
Also $G(r)\ge |g(0)|^2>0$.

Local Lipschitz continuity has a direct uniform estimate.
Fix $0<a<b$. On $[a,b]$, the gradient norm of $|g|^2$ is
$2|g||g'|$, bounded by a constant $C$ on $|z|\le b$.
If $z$ is a maximum point at radius $r\in[a,b]$, compare it with
$(u/r)z$ at radius $u\in[a,b]$. Their distance is $|r-u|$, and the segment
between them lies in that disk. The gradient bound gives
$G(r)\le G(u)+C|r-u|$. Reversing the radii gives
$|G(r)-G(u)|\le C|r-u|$.

Thus $G$ is absolutely continuous on every compact positive interval,
and
\[
 \int_a^bG'(u)\,du=G(b)-G(a)>0.
\]
Monotonicity gives $G'\ge0$ almost everywhere. The set of radii inside
$(a,b)$ where $G$ is differentiable and $G'>0$ therefore has positive
measure. Removing the countable set $Q$ leaves such a radius $r\notin Q$.

At every maximum point $z$ of $g$ on this circle, $g(z)\ne0$.
Differentiating $\theta\mapsto |g(e^{i\theta}z)|^2$ at its maximum gives
angular stationarity. The smooth supporting function
$h_z(u)=|g(uz/r)|^2$ satisfies $h_z(u)\le G(u)$ and $h_z(r)=G(r)$.
Since both functions are differentiable at $r$, their derivatives agree
there. This applies to every maximizing $z$, even when the maximum is tied.
Consequently
\[
 \operatorname{Im}\frac{zg'(z)}{g(z)}=0,\qquad
 \operatorname{Re}\frac{zg'(z)}{g(z)}=\frac{rG'(r)}{2G(r)}>0.
\]
All maximum points therefore have the same real actual value
\[
 A(z)=p=m+\frac{rG'(r)}{2G(r)}>m.
\]
The injective map $z\mapsto(z,\bar z)$ sends them into the explicitly
finite fiber $F(r^2,p)$, because
$B(\bar z)=\overline{A(z)}=p$.
Their cardinality is at most $L=2k$.
The maximum-point sets of $f$ and $g$ coincide, so the same bound holds
for $f$. Choosing the interval $(R,R+1)$ proves Theorem~\ref{thm:main}.

\section{Formalization and verification}\label{sec:lean}

The accompanying Lean~4 project proves the theorem with the original
entire, nonzero, and non-monomial hypotheses. Its final declaration is
\begin{center}
\small\nolinkurl{MaximumModulus.bounded_maxPoints_at_arbitrarily_large_radii}
\end{center}
in \path{MaximumModulus/AllRadii.lean}, line~26. The statement spells out
the set using the norm equality and all actual function-value comparisons.
It has one $B:\N$ before the universal quantifier over $R:\R$, and
explicitly supplies finiteness together with the bound on
\texttt{Set.ncard}. Thus the convention for \texttt{Set.ncard} on an
infinite set cannot make the result vacuous. Independent finiteness on
every positive circle is in \path{MaximumModulus/Finiteness.lean},
line~187. The negative-growth implication is in
\path{MaximumModulus/AllRadii.lean}, line~36.

The principal correspondence between the proof and project modules is
as follows. All filenames in the table are relative to
\path{MaximumModulus/}.
\begin{center}
\small
\begin{tabular}{>{\raggedright\arraybackslash}p{0.35\textwidth}
                >{\raggedright\arraybackslash}p{0.57\textwidth}}
Mathematical step & Lean source modules\\[3pt]
Circle finiteness and logarithmic derivative &
\nolinkurl{Finiteness.lean}, \nolinkurl{LogDerivative.lean}\\[3pt]
Local power coordinates and small fibers &
\nolinkurl{LocalValence.lean}, \nolinkurl{SmallFibers.lean}\\[3pt]
Exact finite covers and distinct-pair counts &
\nolinkurl{FiniteTargetCover.lean}, \nolinkurl{DuplicateSafeCount.lean}\\[3pt]
Ordinary finite values and infinity &
\nolinkurl{HighCountInterior.lean}, \nolinkurl{InfinityInterior.lean}\\[3pt]
Omitted-value capture and continuation &
\nolinkurl{BoundaryCharts.lean}, \nolinkurl{BranchContinuation.lean},
\nolinkurl{OmittedTargetCover.lean}\\[3pt]
Global empty high-count locus & \nolinkurl{HighCountGlobal.lean}\\[3pt]
Countable exceptions and radial envelope &
\nolinkurl{FiniteExceptionalSet.lean}, \nolinkurl{MaximumEnvelope.lean}\\[3pt]
Final statement and negative implication & \nolinkurl{AllRadii.lean}
\end{tabular}
\end{center}

The crucial exact-cover, central-count, and omitted-value assertions are
proved by the declarations
\begin{quote}\small
\nolinkurl{sphereCorrespondence_regular_finite_family},\\
\nolinkurl{central_fiber_count_le_of_total_graph},\\
\nolinkurl{omittedHighCountBranchCover_exists},\\
\nolinkurl{OmittedHighCountBranchCover.actual_closure_branches},\\
\nolinkurl{regular_high_count_closure_omitted_local_projection_isOpen},\\
\nolinkurl{regularLogDerivativeHighCount_empty}.
\end{quote}

The exact toolchain and dependency revisions are
\begin{quote}\small
Lean toolchain: \nolinkurl{leanprover/lean4:v4.35.0-rc4}\\
Lean commit: \nolinkurl{c29b6dda4f7c20e3eeaa717c4e565663c5cfa364}\\
mathlib commit: \nolinkurl{81d17696471311f5e3e2f034236df7a781449f7e}.
\end{quote}
All transitive dependency revisions are pinned in \path{lake-manifest.json}.
Build and verification instructions are in the project \path{README.md};
\texttt{sh scripts/verify.sh} performs the recorded project checks.

A clean project build, direct checks of all project sources, sequential
kernel replays of the 42 project modules relative to their pinned imports,
and explicit target axiom inspections have been executed successfully.
The final theorem and its transitive dependencies use only
\texttt{propext}, \texttt{Classical.choice}, and \texttt{Quot.sound}.
They contain no proof placeholders or unproved custom mathematical axioms.
The requested prize verification workflow was also executed locally on a
frozen audit snapshot. Eight theorem and independent bridge targets passed;
the independently specified final statement and exported solution closure
were accepted by the configured comparison and external checker runs.
These checks concern the formal theorem; they do not constitute a
publication or prize decision.

Full commands, logs, axiom reports, pins, comparison outputs, and isolation
details are preserved in \path{verification/prize-report.md} and its linked
evidence. The audit snapshot commit
\nolinkurl{8b731b51dcb0ed3d7e8c6463b6266c35d8099c85} is a temporary local
verification record, not a published repository revision.
The mathematical exposition here was revised after those checks without
changing the verified Lean source, dependency lock, or build scripts.

\section{Scope and differences from the earlier manuscript}\label{sec:scope}

The result proved here is the cofinal finite bound in
Theorem~\ref{thm:main}, which suffices for the original all-radii question.
It remains compatible with unbounded counts along selected radii.
The earlier manuscript also proposed a locally finite exceptional set of
radii for a $2k$ bound and sharpness examples. Those stronger assertions
are outside the verified scope of this paper and its final Lean theorem.
In particular, local finiteness of the nonregular \emph{target} set $D$
does not by itself imply local finiteness of its projection onto radii.
The argument in Section~\ref{sec:radii} uses countability and a
positive-measure set of radii; it is not an assertion of an almost-everywhere
maximum-point bound.

The continuation proof replaces the earlier proposed propagation of global
irreducible-component degrees by the closure of actual regular high-count
targets. The replacement makes explicit the total-image graph condition,
the exact finite cover, the central distinct-pair inequality, simultaneous
target and parameter shrinking, and the reciprocal coordinate at infinity.
At the omitted value it captures a bounded surviving pair for each target
and derives branch selection from closed and open membership in $H$.
It never assumes that the partial family counts escaping pairs.
The original manuscript has been preserved separately without alteration.

\paragraph{Assistance disclosure.}
OpenAI's ChatGPT and Codex (GPT-6) were used to develop the proof and
prepare the manuscript. The human author takes full responsibility for
the mathematical arguments and the final manuscript.
Codex also assisted with the Lean formalization.
The accompanying verification records describe the checks actually
performed and their scope.

\begin{thebibliography}{9}
\raggedright

\bibitem{HL}
W.~K.~Hayman and E.~F.~Lingham,
\emph{Research Problems in Function Theory},
arXiv:1809.07200v2 (2018), Problem~2.16, printed p.~29.
\url{https://arxiv.org/pdf/1809.07200v2#page=30}.

\bibitem{HP}
F.~Herzog and G.~Piranian,
\emph{The counting function for points of maximum modulus},
in \emph{Entire Functions and Related Parts of Analysis}
(Proc. Sympos. Pure Math., La Jolla, 1966),
American Mathematical Society, Providence, RI, 1968, pp.~240--243.

\bibitem{GP}
Adi Gl\"ucksam and Leticia Pardo-Sim\'on,
\emph{An approximate solution to Erd\H{o}s' maximum modulus points problem},
arXiv:2208.11154v2 (2023), Question~1.1.
\url{https://arxiv.org/html/2208.11154v2#S1}.

\bibitem{PS}
Leticia Pardo-Sim\'on and David J.~Sixsmith,
\emph{A finite-order answer to a problem of Erd\H{o}s on maximum modulus points},
arXiv:2607.09462v1 (2026), introduction.
\url{https://arxiv.org/html/2607.09462v1#S1}.

\end{thebibliography}

\end{document}
